% --------------------------------------------------
% LLNL Proposal Template
% --------------------------------------------------
\documentclass[11pt]{article}
\usepackage[letterpaper, margin=1in]{geometry}
\usepackage{times}

% --------------------------------------------------
% Colors
% --------------------------------------------------
\usepackage[table]{xcolor}
\usepackage{tikz}

\definecolor{secblue}{HTML}{3B5E7F}
\definecolor{subsecblue}{HTML}{4F81BD}
\definecolor{hdrgray}{HTML}{A7B7C7}

\definecolor{boxblue}{HTML}{587590}
\definecolor{boxgray}{HTML}{E0E6E6}

% --------------------------------------------------
% Draft watermark -- uncomment for drafts
% --------------------------------------------------
%\usepackage[colorspec=.9]{draftwatermark}
%\SetWatermarkText{DRAFT}
%\SetWatermarkScale{1}

% --------------------------------------------------
% Extra packages
% --------------------------------------------------
\usepackage[subfigure]{tocloft}
\usepackage{inconsolata}
\usepackage{abstract}
\usepackage{subfigure}
\usepackage{graphicx}
\usepackage{hyperref}

\hypersetup{
  colorlinks,
  citecolor=violet,
  linkcolor=secblue,
  urlcolor=violet,
  breaklinks
}
\usepackage{wrapfig}
\usepackage{url}
\usepackage{breakurl}
\usepackage{pdfpages}
\usepackage{ctable}
\usepackage{multirow}
\usepackage{appendix}
\usepackage{lipsum}
\usepackage{mdframed}
\usepackage{fancyvrb}
\usepackage{tabu}
\usepackage{tabularx}
\usepackage{multirow}
\usepackage{xspace}

\usepackage{pdfpages}

% Sort multiple citations by citation number
\usepackage[sortcites=true]{biblatex}
\addbibresource{strong-foundations-for-sustainability-fy23.bib}

% Space reduction
\usepackage{caption}
\DeclareCaptionFont{subsecblue}{\color{subsecblue}}
\captionsetup{labelfont={small,subsecblue,bf},font=small}

\usepackage{enumitem}
\setlist{noitemsep,topsep=2pt}

% --------------------------------------------------
% Section heading format
% --------------------------------------------------
\usepackage{titlesec}

\titlespacing{\parameterization}{0em}{.4em}{.5em}
\titlespacing{\section}{0em}{.4em}{.1em}
\titlespacing{\subsection}{0em}{.4em}{.1em}
\titlespacing{\subsubsection}{0em}{.4em}{.1em}
\titlespacing{\paragraph}{0em}{.1em}{.5em}

\def\secformat{\sffamily\bfseries}

\titleformat{\section}{\large\secformat\color{secblue}}{\thesection}{1em}{}
\titleformat{\subsection}{\large\secformat\color{subsecblue}}{\thesubsection}{1em}{}
\titleformat{\subsubsection}{\normalsize\secformat\color{subsecblue}}{\thesubsubsection}{1em}{}
\titleformat{\paragraph}[runin]{\small\secformat}{\theparagraph}{1em}{}

\newcommand{\tablesec}[1]{%
        \vspace{.1em}
        \textbf{\color{subsecblue} #1}\newline}

\setlength{\parindent}{0em}
\setlength{\parskip}{.5em}

% --------------------------------------------------
% Make TOC fonts match sections
% --------------------------------------------------
\renewcommand\abstractnamefont{\normalsize\secformat\color{secblue}}
\renewcommand\contentsname{Table of Contents}
\renewcommand\cfttoctitlefont{\small\secformat\color{secblue}}
\renewcommand\cftsecfont{\small\secformat\color{secblue}}
\renewcommand\cftsubsecfont{\small\sffamily\color{subsecblue}}
\renewcommand\cftsubsubsecfont{\small\sffamily\color{subsecblue}}

% --------------------------------------------------
% Use \todo and \ncite to add some red text where
% more text or a citation
% --------------------------------------------------
\newcommand{\todo}[1]{{\bf \color{red} #1}}
\newcommand{\ncite}{{\color{red}[cite]}}

% --------------------------------------------------
% ADD YOUR TITLE HERE (it's used in several places)
% --------------------------------------------------
\newcommand{\mytitle}{My Awesome Proposal}
\newcommand{\foanumber}{DE-FOA-XXXXXXX}
\newcommand{\foaname}{Extreme Paperwork Scalability}

% --------------------------------------------------
% Headers & Geometry
% --------------------------------------------------

% Header for other pages
\usepackage{fancyhdr}

% Header for first page
\fancypagestyle{first}{
        \lhead{\includegraphics[width=2.7in]{images/llnl-logo.png}}
        \rhead{LLNL-PROP-XXXXXX\\
        Program Announcement Number: \foanumber\\
        \foaname}
        \cfoot{}
}


\renewcommand{\headrulewidth}{0pt}
\pagestyle{fancy}
\lhead{
        \begin{tikzpicture}
                \node[white,draw=white,shape=rectangle,fill=hdrgray] at (-8, 12)
                {\sf\footnotesize \hbox to \textwidth{
                        \mytitle
                        \hfill
                        FY23 ASCR Proposal}
                };
        \end{tikzpicture}
}
\chead{}
\rhead{}

% --------------------------------------------------
% Title, Author, Proposal info table
% --------------------------------------------------
\def\lbox{\cellcolor{boxblue}\color{white}\small\sffamily\bfseries}
\def\rbox{\cellcolor{boxgray}\small}
\taburulecolor{white}

%
% Title
\title{\sf\huge\color{secblue}\mytitle}

%
% Author Box is actually all the metadata they want
%
\author{
\def\arraystretch{1.4}
\begin{tabularx}{0.8\textwidth}{r X}
  \lbox Applicant/Institution      & \rbox Lawrence Livermore National Laboratory \\\hline
  \lbox Address                    & \rbox 7000 East Ave L-557, Livermore, CA, 94550 \\\hline
  \hline
  \hline
  \lbox Lead PI Name               & \rbox Your Name \\\hline
  \lbox Telephone                  & \rbox (925) 42X-XXXX \\\hline
  \lbox Email                      & \rbox \href{mailto:you@llnl.gov}{you@llnl.gov} \\\hline
  \hline
  \hline
  \lbox Administrative POC         & \rbox Chung Bothwell \\\hline
  \lbox Telephone                  & \rbox (925) 423-3614 \\\hline
  \lbox Email                      & \rbox \href{mailto:bothwell1@llnl.gov}{bothwell1@llnl.gov} \\\hline
  \hline
  \hline
  \lbox FOA Number                 & \rbox DE-FOA-XXXXXX \\\hline
  \lbox DOE/SC Program Office      & \rbox Office of Advanced Scientific Computing Research \\\hline
  \lbox DOE/SC Technical Contacts  & \rbox Program Manager Name\\\hline
  \hline
  \hline
  \lbox Research Area              & \rbox \foaname\\\hline
\end{tabularx}
}
\date{}

% --------------------------------------------------
% Main document
% --------------------------------------------------
\begin{document}
\pagenumbering{roman}
\setcounter{page}{1}

% --------------------------------------------------
% Project Summary
% --------------------------------------------------
\thispagestyle{first}

{~\vfill~\sf\huge\color{secblue}%
  \hfill \mytitle \hfill
}

\begin{center}\small\def\arraystretch{0.96}
  \begin{tabular}{lll}
    Name 1 & Institution 1 & Role 1 \\
    Name 2 & Institution 2 & Role 2 \\
  \end{tabular}
\end{center}

\paragraph{Project Summary.}
Proposals have always been a lot of work, but in recent years the complexity of
submitting one has grown expoentially. It is projected that by the end of the 2020's,
ASCR proposals will require more paperwork than one can fit in the average HPC
center's tape archive. This proposal will pioneer techniques to allow us to do more
paperwork than ever before, saving valuable machine cycles and using much more
efficient humans to fill in all the appopriate forms.

\vfill

% --------------------------------------------------
% Title Page
% --------------------------------------------------
\newpage
\maketitle

% --------------------------------------------------
% Budget Breakdown Page
% --------------------------------------------------
\newpage
\section*{Summary of Subawards and Project Structure}
\hypertarget{sec:budget-summary}{}
\taburulecolor{black}
\begin{center}
  \begin{tabular}{ |l|p{1.5in}|l|r|r|r| }
    \hline
    \multicolumn{6}{ |c| }{\bf Budget Summary}\\\hline
    \hline
            & {\bf Name} & {\bf Institution}  & {\bf Year 1}  & {\bf Year 2}  & {\bf Total} \\\hline
    Lead PI & Your Name  & Institution 1      & \$XXX,XXX     & \$XXX,XXX     & \$XXX,XXX   \\\hline
    Co-PI   & Their Name & Institution 2      & \$XXX,XXX     & \$XXX,XXX     & \$XXX,XXX   \\\hline
  \end{tabular}
\end{center}

\subsection*{Leadership Structure}

Description of Leadership Structure.

\subsection*{Facilities, Equipment, and Resources}

LLNL Facilities and Equipment are summarized in Appendix~\ref{sec:facilities}.

\subsection*{Mentoring}

Mentoring and outreach efforts for students and early-stage researchers are summarized
in the PIER in Appendix~\ref{sec:pier}.

% --------------------------------------------------
% Table of contents
% Tweak the first baselinestretch below to compress
% it if you need to.
% --------------------------------------------------
\newpage
\renewcommand{\baselinestretch}{1.0}\normalsize
\tableofcontents
\renewcommand{\baselinestretch}{1.0}\normalsize


% --------------------------------------------------
% Narrative
% --------------------------------------------------
\newpage
\pagenumbering{arabic}
\section{Narrative}

Proposals require all kinds of paperwork.

\subsection{Section 1}
We list the various traditional kinds here.

\subsection{Section 2}
New kinds of paperwork are constantely being developed by the DOE. It is hard to keep
up.

\subsection{Section 3}
Our techniques will revolutionize paperwork handling within DOE, saving time for
much-valued machines and leveraging the humans for the hard parts.

\subsection{Conclusions}
We will need to develop even more advanced paperwork techniques in the future. Hopefully
the techniques developed here will get us through this funding round, though.

% --------------------------------------------------
% References
% --------------------------------------------------
\newpage
\appendix % start appendices here

\begingroup
\section{References}
\printbibliography[heading=none]
\endgroup

% --------------------------------------------------
% Letters of Support
% -------------------------------------------------
\newpage
\section{Letters of Support}

Letters of support in alphabetical order.

%\begin{itemize}
%\item \hyperlink{letter-one}{First Letter}
%\item \hyperlink{letter-two}{Second Letter}
%\end{itemize}

%\clearpage
%\hypertarget{letter-one}{}
%\includepdf[pages=-]{letters-of-support/one.pdf}
%\hypertarget{letter-two}{}
%\includepdf[pages=-]{letters-of-support/two.pdf}

% --------------------------------------------------
% Facilities
% -------------------------------------------------
\section{Facilities and Other Resources} \label{sec:facilities}
See the following pages for a description of the Livermore computing facilities
that are available to the project team.

\includepdf[pages=-]{includes/lc-facilities}

% --------------------------------------------------
% Equipment
% -------------------------------------------------
\section{Equipment}

No special equipment is needed for this project.

% --------------------------------------------------
% Data Management
% -------------------------------------------------
\newpage
\section{Data Management Plan}

% Based on CASC-Data-Management-Plan_2022-05-24_3.pdf from
% https://computing-int.llnl.gov/casc/casc-templates

\noindent {\bf Data Types and Sources.}
The primary data artifacts of the proposed research will be implementations of
new algorithms and methodologies in software. Secondary data artifacts will
pertain to the digital research data needed to validate published research
findings: the input files, input data, and metadata (e.g., number of processors,
threads, etc.) used to generate any published data; instructions or scripts that
enable others to reproduce simulations; instructions or scripts required for
data analysis; data analysis results used to produce tables, plots, and charts
in publications; and, as appropriate, raw output data if it is impractical for
others to reproduce the simulations.

\medskip

\noindent {\bf Content and Format.}
Software libraries will be released as source code with basic documentation.
Source for research (non-production) code will be made available with sufficient
documentation to facilitate validation of published research results. Secondary
data artifacts will be accompanied with text or PDF files explaining the
provenance, data format, and way in which the data was used in the published
research.

\medskip

\noindent {\bf Sharing and Preservation.}
\begin{itemize}
\item As allowed by Laboratory policies for the transmittal of data, supported
  software packages will be made available as source code for download from the
  LLNL Software Portal: \url{http://software.llnl.gov}. All software will be
  made available for at least three years from the conclusion of the date of
  publication.
\item Similarly, source code for unsupported software and secondary data
  artifacts will be made available by the PI as communicated within published
  research results. Due diligence will be made to curate the input files and
  scripts required to execute published test cases as well as the resulting data
  and data analysis scripts, if applicable, so that results may be independently
  validated.
\item At the PI's discretion, secondary data artifacts may be made available
  through the Green Data Oasis (GDO), a large data store on the unrestricted
  LLNL network. This facility is intended to facilitate the sharing of
  scientific data with external collaborators. The GDO consists of two Grid
  Interface Front-End Nodes (GrIFENs), a switch, and seven Scalable Storage
  Units (SSUs), which provide a combined capacity of 620TB of RAID storage.
  Solaris zones will be used to create one or more virtual hosts for each
  project, allowing each project to have its own IP addresses, run their own
  services, and do so in such a way as to not interfere with other projects.
  Fair share resource management will be used to prevent one project from using
  resources to the exclusion of other projects.
\item When appropriate, community-standard test data and problems that are made
  publicly available by third parties will be used. Responsibility for
  maintaining these standard tests will remain with these third parties.
\item As allowed by Laboratory policies for the transmittal of data, data will
  be released under a standard open source license.
\item A concerted effort will be made to use non-proprietary software to build,
  install, run, or analyze the results of software developed and released by
  this project. However, proprietary tools and software will be used when no
  suitable non-proprietary option exists, and the use of such proprietary tools
  will significantly advance the goals of the research.
\item During the duration of the project, all data will be managed under a
  version control system, including tagged versions of the software used to
  produce specific research results. At the conclusion of the project, the
  repository will be retained for no less than three years. An archive of the
  repository will be made on HPSS archival storage.
\item In the event the PI leaves the Laboratory, the institution will retain a
  copy of all data, and the PI will designate a successor to handle subsequent
  requests for data.
\item Should LLNL adopt an institutional plan for data management consistent
  with DOE data management policies, this institutional plan will supersede the
  plan described here.
\end{itemize}

\noindent {\bf Protection.}
In compliance with Federal, State, and local statues and regulations, Lawrence
Livermore National Laboratory has established standards and policies regarding
the transmittal, storage, and destruction of data to protect confidentiality,
personal privacy, Personally Identifiable Information, and U.S. national,
homeland, and economic security; recognize proprietary interests, business
confidential information, and intellectual property rights; and avoid
significant negative impact on innovation and U.S. competitiveness. The PI will
conduct all work in accordance to these standards with support from LLNL's
Office of Classification and Export Control and the Innovation and Partnerships
Office.

\medskip

\noindent {\bf Rationale.}
The primary products for research in this area (for the purposes of this data
management plan) are algorithms, test cases, and/or the input, output, and
performance data as described and documented in peer-reviewed papers and
implemented in software. Thus, the main focus of the proposed data management
plan is on providing access to the software produced and the necessary data to
reproduce published results. Making correct implementations and test data
available to the research community will allow others to validate the results
and provide common problems for comparison of alternative approaches. In the
broader scientific computing community, it also will allow for more rapid
adoption of new techniques.

% --------------------------------------------------
% PIER Plan
% -------------------------------------------------
\newpage
\section{Plan for Promoting Inclusive and Equitable Research (PIER)}
\label{sec:pier}

Your PIER Plan Here.

% --------------------------------------------------
% Biographical Sketches
% --------------------------------------------------
\newpage
\section{Biographical Sketches}

Biographical sketches of senior/key personnel in alphabetical order.

%\begin{itemize}
%\item \hyperlink{bio-one}{Your Name}
%\item \hyperlink{bio-two}{Your Co-PI}
%\end{itemize}

%\clearpage
%\hypertarget{bio-one}{}
%\includepdf[pages=-]{bios/bio-you.pdf}
%\hypertarget{bio-two}{}
%\includepdf[pages=-]{bios/bio-copi.pdf}

% --------------------------------------------------
% Current and Pending Support
% -------------------------------------------------
\newpage
\section{Current and Pending Support}

Current and pending support for senior/key personnel in alphabetical order.

%\begin{itemize}
%\item \hyperlink{cp-one}{Your Name}
%\item \hyperlink{cp-two}{Your Co-PI}
%\end{itemize}

%\clearpage
%\hypertarget{cp-one}{}
%\includepdf[pages=-]{cp/cp-you.pdf}
%\hypertarget{cp-two}{}
%\includepdf[pages=-]{cp/cp-copi.pdf}


% --------------------------------------------------
% Budget
% -------------------------------------------------

\newpage
\section{Lead Institution Budget}

The following pages have the budget and budget justification for the lead institution:
Lawrence Livermore National Laboratory. Budgets for subawards at other organizations are
in the next section. Refer to the budget summary in the
\hyperlink{sec:budget-summary}{Summary of Subawards and Project Structure} section for
the budget breakdown.

As the lead institution, LLNL will ...

%\includepdf[pages=-]{budgets/budget.pdf}
%\includepdf[pages=-]{budgets/budget-justification.pdf}

% --------------------------------------------------
% Subawards
% -------------------------------------------------

\newpage
\section{Subaward Budgets}

Subawards in alphabetical order. Each subaward includes a statement of work, budget, and
budget justification.

%\begin{itemize}
%\item \hyperlink{sub-one}{Institution 2}
%\item \hyperlink{sub-two}{Institution 3}
%\end{itemize}

%\clearpage
%\hypertarget{sub-one}{}
%\includepdf[pages=-]{budgets/one/one-sow.pdf}
%\includepdf[pages=-]{budgets/one/one-budget.pdf}
%\includepdf[pages=-]{budgets/one/one-budget-justification.pdf}

%\hypertarget{sub-two}{}
%\includepdf[pages=-]{budgets/two/two-sow.pdf}
%\includepdf[pages=-]{budgets/two/two-budget.pdf}
%\includepdf[pages=-]{budgets/two/two-budget-justification.pdf}


\end{document}
